\section{Conclusion}
\label{sec:conclusion}

We show that the concept ``X is a Pok\'{e}mon'' is encoded as a linear direction in GPT-2-XL's residual stream, concentrated in middle layers (10--32). A mass-mean probe achieves 88.4\% classification accuracy on held-out entities, outperforming both logistic regression and random baselines. By projecting non-Pok\'{e}mon entities onto this direction, we discover a continuous, interpretable spectrum of Pok\'{e}mon quality: game creatures from similar franchises and invented monster names score highest, Digimon rank in the middle, and common objects are maximally distant.

The key takeaway is that linear representations in language models extend to fine-grained, culturally constructed categories---not just broad semantic properties---and that these representations encode graded similarity rather than binary membership.

Future work should validate these findings with causal intervention experiments, replicate across larger models, and investigate whether sub-categories (Pok\'{e}mon types, generations) form additional structure within the Pok\'{e}mon subspace. Sparse autoencoder decomposition may reveal whether a discrete ``Pok\'{e}mon feature'' exists alongside the linear direction we identify.